\begin{itemize}
\item \textbf{NORMAL.} \emph{Numérica:} Newton (6 iteraciones) y secante (8) alcanzan la raíz de referencia hasta el redondeo; la bisección necesita 34 iteraciones porque reduce el intervalo a la mitad en cada paso. En el sistema ($\kappa_{2}$ = 1.373), Gauss-Seidel requiere 8 iteraciones frente a 13 de Jacobi ($\rho$\_GS $\approx$ $\rho$\_J$^{2}$) y el gradiente conjugado 5; LU es exacto hasta el redondeo. \emph{Problema:} El indicador alcanza el umbral de 40 \% a los 23.589 meses. Un retraso de 3 meses lleva el desgaste a 45.411 \% (5.41 \% respecto del umbral; negativo indica que la curva ya volvió a bajar). La fecha es más sensible al umbral y a la tasa $\alpha$ cuanto menor es la pendiente D$'$(t*) = 1.744 \%/mes.
\item \textbf{ESTRÉS.} \emph{Numérica:} Newton (6 iteraciones) y secante (7) alcanzan la raíz de referencia hasta el redondeo; la bisección necesita 44 iteraciones porque reduce el intervalo a la mitad en cada paso. En el sistema ($\kappa_{2}$ = 20.97), Gauss-Seidel requiere 147 iteraciones frente a 285 de Jacobi ($\rho$\_GS $\approx$ $\rho$\_J$^{2}$) y el gradiente conjugado 22; LU es exacto hasta el redondeo. \emph{Problema:} El indicador alcanza el umbral de 40 \% a los 5.3518 meses. Un retraso de 3 meses lleva el desgaste a 52.167 \% (12.2 \% respecto del umbral; negativo indica que la curva ya volvió a bajar). La fecha es más sensible al umbral y a la tasa $\alpha$ cuanto menor es la pendiente D$'$(t*) = 3.606 \%/mes.
\item \textbf{MAL COND..} \emph{Numérica:} Las tres raíces convergen (bisección 34, Newton 17, secante 24 iteraciones), pero D$'$(t*) = 4.472$\times10^{-4}$ hace que el problema sea mal condicionado: $\kappa$ = 3727.47 y una perturbación relativa de 10$^{-9}$ en D\_u se amplifica unas $\kappa$ veces, mientras que +10$^{-7}$ elimina la raíz (el umbral deja de alcanzarse). En el sistema de Hilbert ($\kappa_{2}$ = 1.526$\times10^{10}$) Jacobi diverge ($\rho$ = 6.04 > 1), Gauss-Seidel no termina, y LU y gradiente conjugado dejan residuos mínimos con errores mucho mayores: mal condicionamiento y falta de convergencia son fenómenos distintos. \emph{Problema:} El indicador alcanza el umbral de 40 \% a los 23.996 meses. Un retraso de 3 meses lleva el desgaste a 39.551 \% (-0.449 \% respecto del umbral; negativo indica que la curva ya volvió a bajar). Como la curva casi toca el umbral, la fecha de inspección depende de forma extrema del umbral elegido y de los coeficientes: no debe planificarse con un solo valor, sino con un margen o con un umbral menor.
\end{itemize}
